\documentclass[a4paper, 12pt]{article}

\usepackage{utils}
\renewcommand*{\today}{27 March 2026}

\binoppenalty=10000
\relpenalty=10000

\begin{document}

\hotbox{Algèbre bilinéaire}{CM 9}{\today}

\begin{proposition}{}{}
    Soit $(E, \langle \cdot, \cdot \rangle)$ un espace vectorial euclidien orienté de dimension finie~$n$.

    Soient $x_1, \dots, x_{n-1} \in E$, alors on a

    \begin{enumerate}[label=(\alph*)]
        \item La famille $(x_1, \dots, x_{n-1})$ est liée si et seulement si $x_1 \wedge \dots \wedge x_{n-1} = 0$.
        
        \item Le produit vectoriel est l'unique application linéaire $f : E^{n-1} \to E$ telle que
        \begin{enumerate}[label=(\roman*)]
            \raggedright
            \item $\forall x_1, \dots, x_{n-1} \in E, f(x_1, \dots, x_{n-1}) \in Vect(x_1, \dots, x_{n-1})^\perp$
            \item $(x_1, \dots, x_{n-1})$ libre implique que $(x_1, \dots, x_{n-1}, f(x_1, \dots, x_{n-1}))$ est une base directe de $E$.
            \item $\forall x \in E^{n-1}, \norm{f(x)}^2 = gram(x)$
        \end{enumerate}
    \end{enumerate}
\end{proposition}

\begin{demonstration}
    \begin{hotwarn}
        À revoir
    \end{hotwarn}
\end{demonstration}

\begin{methode}
    Ainsi, étant donné $(x_1, \dots, x_{n-1}) \in E^{n-1}$ et une application linéaire $f : E^{n-1} \to E$.

    On peut vérifier que $f$ est le produit vectoriel en vérifiant les trois propriétés suivantes :
    \begin{enumerate}[label=(\roman*)]
        \item $f(x_1, \dots, x_{n-1})$ est orthogonal à tous les vecteurs de la famille $(x_1, \dots, x_{n-1})$.
        \item $(x_1, \dots, x_{n-1}, f(x_1, \dots, x_{n-1}))$ est une base directe de $E$.
        \item $gram(x_1, \dots, x_{n-1}) = \norm{f(x_1, \dots, x_{n-1})}^2$.
    \end{enumerate}
\end{methode}

\end{document}